\documentclass[10pt,a4paper]{amsart}
\usepackage[margin=19mm]{geometry}
\usepackage{amsmath,amssymb,mathtools}
\usepackage[scr=rsfs,cal=euler]{mathalfa}
\usepackage{xfrac}
\usepackage[unicode,hidelinks,bookmarks=false]{hyperref}
\newcommand{\ie}{i.e.\ }
\newcommand{\cf}{cf.\ }
\newcommand{\resp}{respectively\ }
\newcommand{\Subsection}{\subsection}
\providecommand{\nobreakdash}{\nobreak\hbox{-}}
\setlength{\emergencystretch}{2em}
\multlinegap0pt
\usepackage{xcolor}
\usepackage{cancel}
\usepackage{amssymb}
\usepackage{mathtools}
\usepackage{float}
\usepackage[shortlabels]{enumitem}
\newtheorem{thm}{Theorem}
\newtheorem{prop}{Proposition}
\newtheorem{lem}{Lemma}
\newcommand{\E}{\mathbb E}
\newcommand{\Gini}{\mathcal D}
\newcommand{\Unif}{\operatorname{Unif}}
\newcommand{\Var}{\operatorname{Var}}
\title{The Zaporozhets–Tarasov Conjecture on Mean Distances}
\author{Eric Shen}
\date{Source: arXiv:2608.06470v3 (CC BY 4.0)}
\begin{document}
\maketitle

\noindent\emph{Source and licence.} The Zaporozhets–Tarasov Conjecture on Mean Distances. Version arXiv:2608.06470v3 (CC BY 4.0), distributed by arXiv under the Creative Commons Attribution 4.0 International licence, http://creativecommons.org/licenses/by/4.0/, as recorded in the arXiv licence metadata for this exact version and shown on the abstract page https://arxiv.org/abs/2608.06470v3. Full licence notice: \texttt{licenses/arxiv-2608.06470-CC-BY-4.0.txt}.\par\smallskip

\noindent\textit{Editorial scope.} Counted unit: the article's lem lem:upper\_bound (source0.tex lines 266--269). The article's complete original proof of it is delivered with it (source0.tex lines 280--483, one \texttt{proof} environment of 204 lines, which the article places away from the statement). No mathematics is written, repaired or re-derived: the statement and the proof are the article's own lines, in the article's own notation, and no retained line is substituted or re-flowed.  Retained uncounted as the source-local context the proof needs: lines 137--161; lines 192--203; lines 219--225.  Nothing was omitted from any retained span, so the package carries no omission record.  No retained line contains a citation, so the wrapper carries no bibliography and no bibliography entry.  No retained line contains a figure environment, an image-inclusion command or drawing code, so no image is delivered and the package declares no asset dependency.  Editorial formatting only, in addition: an AMS \texttt{amsart} preamble that restates the article's own declarations and macros used by the retained lines, namely the environments \texttt{thm}, \texttt{prop}, \texttt{lem} on separate counters, together with the standard packages those lines need.  The retained environments print in this document as Theorem 1, Proposition 1, Lemma 1 in order of appearance, whereas the article prints the same items with the numbering of its own full document.  The complete native source of the article as distributed in the arXiv e-print for this exact version is bundled unchanged as \texttt{sources/207000/source0.tex}.
\begin{thm}
    \label{thm:endpoint}
Let $\varphi$ be a nonincreasing probability density and $\psi$ a nondecreasing probability density on $[0,1]$. Let
$$
 \Phi(x)=\int_0^x\varphi(s)\,ds,
 \qquad
 \Psi(x)=\int_0^x\psi(s)\,ds.
$$
Put
$$
 H=\Phi-\Psi,
 \qquad
 A=\int_0^1H(x)\,dx,
$$
and assume $A>0$. Define
$$
 p(x)=\frac{H(x)}A,
 \qquad
 q(x)=\frac{\varphi(x)+\psi(x)}2.
$$
Then
$$
 \Gini(q)\ge \Gini(p).
$$
\end{thm}

\begin{prop} \label{prop:monotone_factoriz}
Let $\varphi$ be a probability density on $[0,1]$ admitting a nonincreasing representative, which we also denote by $\varphi$. There is a random variable $T\in[0,1]$ such that, for $U\sim\Unif[0,1]$ independent of $T$, the variable
$$
 X=TU
$$
has density $\varphi$. Likewise, for every nondecreasing density $\psi$, there is $S\in[0,1]$ such that
$$
 Y=1-SV,
 \qquad V\sim\Unif[0,1],
$$
has density $\psi$.
\end{prop}

\begin{lem}\label{lem:low_bound_enpdpoint}
    In the notation of Theorem~\ref{thm:endpoint} and Proposition~\ref{prop:monotone_factoriz}, denote 
    $$ A=\int_0^1(\Phi-\Psi)
   =\E Y-\E X
   =1-\frac{\E T+\E S}{2}.$$
   Then $\Gini(q) \geq \frac{1}{3} B := \frac{1}{3}(1-A+\frac{3}{2}A^2 + \Var(T) + \Var(S)).$
\end{lem}

\begin{lem} \label{lem:upper_bound} Let $p$ be as in Theorem~\ref{thm:endpoint}, let $Z$ have density $p$, and assume $B$ as in Lemma~\ref{lem:low_bound_enpdpoint}. Then
$$ 12\Var Z\le B^2.$$

\end{lem}

\begin{proof}[Proof of Lemma~\ref{lem:upper_bound}]
This part is rather technical, though quite elementary. For every integer $m\ge0$, Fubini's theorem gives
$$
 \E Z^m
 =\frac{\E Y^{m+1}-\E X^{m+1}}{(m+1)A}.
$$
Indeed, 
$$
 \int_0^1x^m\Phi(x)\,dx = \int_0^1x^m\E\mathbf 1_{\{X\leq x\}}\,dx
 =\E\int_X^1x^m\,dx
 =\frac{1-\E X^{m+1}}{m+1},
$$


and since $p=(\Phi-\Psi)/A$, the analogous identity for $\Psi$ gives the claim after subtraction. What follows is mostly algebraic identities involving $T$ and $S$.

 Recall that $X = TU$, and $Y = 1 -SV$, and introduce
$$
 z=\frac{\E T-\E S}{2},\quad
 e=\E[T(1-T)]+\E[S(1-S)],\quad
 k=\E[T(1-T)]-\E[S(1-S)].
$$


Using the general moment formula with $m=1$, and using
$$
 \E X^2=\frac{\E[T^2]}{3},
 \qquad
 \E Y^2
 =1-\E S+\frac{\E[S^2]}{3},
$$
we obtain
$$
 \E Z
 =
 \frac{
  1-\E S+\frac13\bigl(\E[S^2]-\E[T^2]\bigr)
 }{2A} =\frac12+\frac{z+k}{6A}.
$$

Similarly, using the moment formula with $m=2$,
$$
 \E X^3=\frac{\E[T^3]}4,
$$
while
$$
 \E Y^3
 =
 \E(1-SV)^3
 =
 1-\frac32\E S+\E[S^2]-\frac14\E[S^3].
$$
Hence
$$
 \E Z^2
 =
 \frac{
  1-\frac32\E S+\E[S^2]
  -\frac14\bigl(\E[S^3]+\E[T^3]\bigr)
 }{3A}.
$$
The numerator can be regrouped as
$$
 1-\frac32\E S+\E[S^2]
  -\frac14\bigl(\E[S^3]+\E[T^3]\bigr) =
 A+\frac{z+k}{2}
 -\frac14\left(
   \E[T(1-T)^2]+\E[S(1-S)^2]
 \right). $$
Consequently,
$$
 \E Z^2
 =
 \frac13+\frac{z+k}{6A}
 -\frac{
   \E[T(1-T)^2]+\E[S(1-S)^2]
 }{12A}.
$$

Subtracting $(\E Z)^2$ gives
$$ 12\Var Z =
 12\E Z^2-12(\E Z)^2 = 1 -\frac{\E[T(1-T)^2]+\E[S(1-S)^2]}{A} -\frac{(z+k)^2}{3A^2}.$$

We now estimate the two expectations which occur with a minus sign.
By Cauchy--Schwarz,
$$
 \bigl(\E[T(1-T)]\bigr)^2
 \le
 \E T\;\E[T(1-T)^2].
$$
Since
$$
 \E[T(1-T)]=\frac{e+k}{2},
 \qquad
 \E T=1-A+z,
$$
this gives
$$
 \E[T(1-T)^2]
 \ge
 \frac{(e+k)^2}{4(1-A+z)}.
$$
Likewise,
$$
 \E[S(1-S)^2]
 \ge
 \frac{(e-k)^2}{4(1-A-z)}.
$$
Substitution into the preceding formula yields
$$
 12\Var Z\le 1-\Lambda,
$$
where
$$
 \Lambda=
 \frac{(e+k)^2}{4A(1-A+z)}
 +\frac{(e-k)^2}{4A(1-A-z)}
 +\frac{(z+k)^2}{3A^2}.
$$


We claim that
$$
 \Lambda\ge2(1-B),
$$
where $B=1-A+\frac{3}{2}A^2 + \Var(T) + \Var(S) = 1+A-\frac{1}{2}A^2-2z^2-e$.
\par Once this is proved, the result follows immediately. Indeed, $B\ge0$ from its original definition, and therefore
$$
 12\Var Z\le1-\Lambda\le2B-1\le B^2,
$$

It remains to prove the claim. Since
$\E T=1-A+z$ and $\E S=1-A-z$ lie in $[0,1]$,
$$
 0\le z^2\le\min\{A^2,(1-A)^2\}.
$$
Using the latter formula for $B$, we obtain

$$ \Lambda-2(1-B)
 ={}\frac{(e+k)^2}{4A(1-A+z)}
 +\frac{(e-k)^2}{4A(1-A-z)}
 +\frac{(z+k)^2}{3A^2}
 +2A-A^2-4z^2-2e.$$

This is a convex quadratic polynomial in $k$ and $e$, which we estimate from below in a standard way. Recall $1-A+z = \E T$, $1-A-z=\E S$ and, to simplify the analysis, introduce new variables $e=x+y$, $k=x-y$. Then
$$\Lambda-2(1-B) = F(x,y) = \frac{x^2}{A \E T} + \frac{y^2}{A \E S} + \frac{(z+x-y)^2}{3A^2} + 2A-A^2-4z^2-2x-2y.$$
Minimizing that in $x,y$ gives
$$F(x,y) \geq \frac{A^3(A+2)-(4A-1)z^2}{A(A+2)}.$$

% Derivation of the lower bound for F.
%
% Set $$d=z+x-y.$$
%
% At the critical point,
% $$\frac{x}{A\,\E T}+\frac{d}{3A^2}=1,\qquad\frac{y}{A\,\E S}-\frac{d}{3A^2}=1.$$
%
% Hence
% $$x=A\,\E T-\frac{\E T}{3A}d,\qquad  y=A\,\E S+\frac{\E S}{3A}d.$$
%
% Since
% $$\E T-\E S=2z,\qquad\E T+\E S=2(1-A),$$
% we obtain
% $$d=z+x-y=(1+2A)z-\frac{2(1-A)}{3A}d.$$

% Therefore
% $$d=\frac{3A(2A+1)}{A+2}\,z.$$

% To evaluate $F$ at the critical point without substituting
% $x$ and $y$ directly, multiply the two critical-point equations
% by $x$ and $y$, respectively, and add. This gives
% $$\frac{x^2}{A\,\E T}+\frac{y^2}{A\,\E S}=x+y-\frac{d(d-z)}{3A^2}.$$

% Consequently,
% $$F_{\min} =-(x+y)+\frac{zd}{3A^2} +2A-A^2-4z^2.$$
% $$ x+y  = A(\E T+\E S)  +\frac{\E S-\E T}{3A}d  = 2A(1-A)-\frac{2zd}{3A}. $$
% Thus  $$ F_{\min}  = A^2-4z^2+\frac{(2A+1)zd}{3A^2}.$$
% Substituting the expression for $d$ yields $$
% F_{\min} =A^2-4z^2 +\frac{(2A+1)^2}{A(A+2)}z^2 =\frac{A^3(A+2)-(4A-1)z^2}{A(A+2)}.$$



Thus it only remains to verify that the numerator is nonnegative.






If $0<A\le\frac14$, then $4A-1\le0$, so $$A^3(A+2)-(4A-1)z^2\ge A^3(A+2) \geq 0.$$ If
$\frac14\le A\le\frac12$, then $z^2\le A^2$ and
$$
 A^3(A+2)-(4A-1)z^2\ge A^3(A+2)-(4A-1)A^2
 =A^2(1-A)^2\ge0.
$$
Finally, if $\frac12\le A<1$, then $z^2\le(1-A)^2$ and
$$
 \begin{aligned}
 A^3(A+2)-(4A-1)z^2&\ge A^3(A+2)-(4A-1)(1-A)^2\\
  &=(A^2-A)^2+(4A-1)(2A-1)\ge0.
 \end{aligned}
$$
Therefore $\Lambda-2(1-B)\ge0$, and hence $12\Var Z\le B^2$.

\end{proof}

\end{document}
